% Public good accumulation --- plan \S3.3, "\S7".
% Sources: 10/04:1-40.  10/04:41-65 is parked as
% _parked/public_good_closed_group_notes.tex.
% Rename: birth rate \gamma(t) -> \lambda(t) = \lambda_0 + \alpha P(t)
% (plan \S7.2), and the base rate \beta -> \lambda_0 (\beta is the speciation
% rate and nothing else).
% The epistemic hedge of 10/04:4 is carried over near-verbatim, reworded out
% of the first person; see PHASE_A_REPORT.md item 2.

\section{Public good accumulation}
\label{var:sec:publicgood}

\subsection{Where the idea comes from}
\label{var:sec:publicgoodframing}

The last two models concern \yp{}, and this one borrows its mechanism from
work on antimicrobial resistance. That work considers populations of two
types, one resistant to a drug and one susceptible, in which the resistant
type produces a substance and exports it into the surrounding medium. Because
the substance protects every cell it reaches and not only the ones that made
it, it is a \emph{public good}: the susceptible type benefits from the
investment of the resistant one, and the resulting dynamics are not those of a
straightforward competitive exclusion \cite{asker2023coexistence}.

Something similar may hold between \yp{} and the macrophage. \yp{} has
evolved resistance to the toxins presented within macrophages, and were that
resistance executed through a substance of this kind, confinement inside the
host cell would concentrate the substance rather than disperse it --- an
advantage available inside a cell and not outside one. This is purely
speculative and may not at all be the case, but even if it is not, the
dynamics may be interesting in and of themselves. Taken as a hypothesis it
would further support the argument of \cref{path:sec:ypestis} that the
within-macrophage environment carries a replicative advantage, and would
supply a mechanism for that advantage rather than only an inference from
behaviour.

\begin{center}
\small
\captionof{table}{Reading of the public-good model in an intracellular
\yp{} setting. The last three rows are the speculation; the first four are
the process as written.}
\label{var:tab:pg-map}
\begin{tabular}{@{}p{0.32\textwidth}p{0.58\textwidth}@{}}
\toprule
Model component & Biological meaning under this reading \\
\midrule
$X_t$ & intracellular \yp{} count in one macrophage \\
$P(t)$ & concentration of an exported protective substance \\
$g$, $c$ & per-bacterium production, and clearance from the enclosed volume \\
$\lambda(t)=\lambda_0+\alpha P(t)$ & replication raised by the accumulated good \\
$\mu$ & killing by macrophage toxic products, treated as constant \\
Enclosure & concentrates $P$ by preventing diffusion away \\
Public-good mechanism & resistance executed through an exported substance;
  not established for this organism \\
\bottomrule
\end{tabular}
\end{center}

\begin{figure}[H]
\centering
\includegraphics[width=0.78\textwidth]{fig10_public_good/fig10}
\caption{\yp{} replicating inside a macrophage, with public good $P(t)$
accumulating in the enclosed volume. The enclosure stops the good diffusing
away, so the same production rate $gX_t$ buys a higher concentration inside
a cell than outside; clearance $c$ is the competing loss. Replication
responds at $\lambda(t)=\lambda_0+\alpha P(t)$.}
\label{var:fig:board}
\end{figure}

\subsection{The model}
\label{var:sec:publicgoodmodel}

Let $X_t$ count \yp{} within one macrophage and $P(t)$ the concentration of
public good in it. The good is produced in proportion to the bacterial count
and cleared at a constant rate, and its presence raises the replication rate:
\begin{align}
\ddt{P}&=g\,X_t-c,\notag\\
\pr{\Delta X_t=1\mid X_t}&=\lambda(t)\,X_t\,\Delta t+o(\Delta t),
\qquad
\lambda(t)=\lambda_0+\alpha P(t),
\label{var:eq:publicgood}\\
\pr{\Delta X_t=-1\mid X_t}&=\mu\,X_t\,\Delta t+o(\Delta t),\notag
\end{align}
with $g$ the production rate per bacterium, $c$ the clearance rate,
$\lambda_0$ the base birth rate and $\alpha$ the public-good advantage. The
effect of the host cell's toxic products is folded into the constant death
rate $\mu$, which ignores the depletion of those products that \ChPath{}
treats explicitly; the object of interest here is the public good rather than
the toxin.

The feedback is the one this chapter is about, running for once in the
favourable direction: a larger population makes more of the good, which
raises the birth rate, which makes the population larger still. What stops
that being unconditional is the clearance term $-c$. While $X_t<c/g$ the good
is cleared faster than it is made and $P$ falls, so a small founding cohort
may never start the loop at all, and the regime worth attention is therefore
$\mu>\lambda_0$, where the process is subcritical on its own and turns
supercritical only once enough public good has accumulated. Survival there is
not a matter of luck alone but of initial numbers.

\subsection{The minimum founding cohort}
\label{var:sec:kstar}

That observation can be made into a well-posed question, and it is the best
one this model has. Suppose the process embedded in a host-to-host sequence:
if the infection within a macrophage survives to late time the cell dies and
delivers its load into two more cells, and if not, the lineage ends. This is
a Galton--Watson process whose division probability is set by the within-cell
dynamics above, and it is supercritical exactly when the probability of
within-cell survival exceeds one half. The quantity wanted is the smallest
founding cohort achieving that:
\begin{equation}
K^\ast=\inf\left\{K:\ \lim_{t\to\infty}\pr{X_t>0\mid X_0=K}>\tfrac12\right\}.
\label{var:eq:kstar}
\end{equation}
In the regime $\mu>\lambda_0$ the limit is zero for every $K$ once the public
good is removed, so $K^\ast$ is finite only on account of it and measures
exactly what the public good is worth: it is the number of bacteria that must
arrive together for the mechanism to pay for itself.

\Cref{var:fig:kstar} computes it. At
$(\lambda_0,\alpha,g,c)=(0.6,0.25,0.15,0.45)$, so that $P$ falls whenever
$X_t<c/g=3$, the survival probability as a function of $K$ is a steep
sigmoid through one half, and $K^\ast$ itself rises with the excess death
rate $\mu-\lambda_0$. The dashed line in panel~(b) is $c/g$: $K^\ast$ sits
above it, as it must, because a founding cohort smaller than the clearance
threshold cannot start the loop. At $\mu-\lambda_0=1.2$ no $K\le 28$ crossed
one half in $600$ paths, which is a statement about the parameter point
rather than a failure of the definition.

Two things about \cref{var:eq:kstar} are worth stating plainly. It is not a
threshold in the sense of \cref{path:cor:threshold}, which is a statement
about a budget and is available in closed form; this one is defined by a
limit of a survival probability, and there is no reason to expect it in
closed form. The two-offspring assumption is also doing work, though less
than it appears to: a general offspring law with mean $m$ would replace the
one-half in \cref{var:eq:kstar} by $1/m$, which moves $K^\ast$ without
changing what it measures.

\begin{figure}[!htbp]
\centering
\includegraphics[width=0.96\textwidth]{fig14_kstar/fig14}
\caption{The minimum founding cohort of \cref{var:eq:kstar}, computed by
Gillespie simulation of \cref{var:eq:publicgood} with $P$ floored at zero
($600$ paths per point).
Panel~(a): estimated late-time survival against $K$ at $\mu-\lambda_0=0.5$,
with the one-half line that defines $K^\ast$.
Panel~(b): $K^\ast$ against excess death; the dashed line is the clearance
threshold $c/g=3$.
Parameters $(\lambda_0,\alpha,g,c)=(0.6,0.25,0.15,0.45)$.}
\label{var:fig:kstar}
\end{figure}

\openqs{
\item \textbf{Death rate rather than birth rate.} It may well make more sense
for the public good to lower the death rate than to raise the birth rate,
since what it protects against is killing. That is a different model ---
$\mu(t)=\mu_0-\alpha P(t)$ with a floor at zero, which is not a linear
birth--death process in the same way --- and it has yet to be written down.
\item \textbf{Whether the enclosure argument survives contact.} The whole
motivation rests on confinement concentrating the good, and making that
quantitative wants a volume and a diffusion or export rate, neither of which
appears in \cref{var:eq:publicgood}. As written the model assumes the
concentrating effect rather than deriving it.
\item \textbf{The scaling of $K^\ast$.} \Cref{var:fig:kstar} is one slice of
parameter space. How $K^\ast$ scales with $\alpha$ and with $c/g$, and
whether the rise with $\mu-\lambda_0$ remains roughly linear, is open.
}

\FloatBarrier
